% 中文文章模板：XeLaTeX + ctexart，字体集 fandol（离线字体包内置）。
% 结构：标题信息 → 摘要/关键词 → 引言 → 相关工作 → 方法 → 实验与结果
%       → 结论 → 参考文献（BibTeX，GB/T 7714—2015 顺序编码制）。
% 引用：条目写入 references.bib，正文用 \cite 命令加条目键引用；
%       只添加真实、可核验的文献。
\documentclass[UTF8,fontset=fandol,a4paper]{ctexart}

\usepackage[margin=2.5cm]{geometry}
\usepackage{amsmath,amssymb,amsthm}
\usepackage{booktabs}
\usepackage{graphicx}
% 参考文献采用 GB/T 7714—2015（gbt7714 宏包 + gbt7714-numerical 样式）：
% 这是国内期刊与学位论文通行的国家标准，能正确处理中英文混排条目
% （作者超过三人时自动写"等"或"et al."），比 plain/unsrt 更适合中文写作。
% 默认上标引用；需要方括号正文引用 [1] 时取消下一行注释。
\usepackage{gbt7714}
% \citestyle{numbers}
\usepackage{hyperref}
\hypersetup{colorlinks=true, linkcolor=blue, citecolor=blue, urlcolor=blue}

% 定理类环境（中文名称）；证明环境 proof 由 amsthm 提供，ctex 已将其标题设为"证明"。
\theoremstyle{plain}
\newtheorem{theorem}{定理}[section]
\newtheorem{lemma}[theorem]{引理}
\theoremstyle{definition}
\newtheorem{definition}[theorem]{定义}

\title{文章标题}
\author{作者姓名\\单位名称}
% \today 显示该版本内容首次构建的日期（构建时钟按快照固定，保证可复现）；投稿时可直接写出日期或留空。
\date{\today}

\begin{document}
\maketitle

\begin{abstract}
% 用 200～300 字概括研究问题、方法、主要结果与结论。
在此填写摘要。

\medskip
\noindent\textbf{关键词：}关键词一；关键词二；关键词三
\end{abstract}

\section{引言}
% 说明研究背景与动机，界定研究问题，概述本文贡献，并交代全文结构。

\section{相关工作}
% 按主题归纳已有研究并引用文献，指出本文与它们的区别。

\section{方法}
% 给出问题的形式化描述、核心定义与主要结论。以下环境仅示范写法，请替换为实际内容。
\begin{definition}\label{def:main}
在此给出定义。
\end{definition}

\begin{lemma}\label{lem:main}
在此陈述引理。
\end{lemma}

\begin{theorem}\label{thm:main}
在此陈述定理。
\end{theorem}

\begin{proof}
在此给出证明，可引用引理~\ref{lem:main} 与式~\eqref{eq:main}。
\end{proof}

编号公式示例（替换为实际公式）：
\begin{equation}\label{eq:main}
  f(x) = \sum_{i=1}^{n} w_i\, g_i(x) .
\end{equation}

\section{实验与结果}
% 介绍数据、实验设置与评价指标。表中数据必须来自真实实验记录。
\begin{table}[htbp]
  \centering
  \caption{实验结果（表头为示例，数据待填）}
  \label{tab:results}
  \begin{tabular}{lcc}
    \toprule
    方法     & 指标一 & 指标二 \\
    \midrule
    基线方法 & --     & --     \\
    本文方法 & --     & --     \\
    \bottomrule
  \end{tabular}
\end{table}

% 插图：图片放入 figures/ 目录后，用 figure 环境加 \includegraphics 命令插入，
% 并配 \caption 与 \label，正文用 \ref 引用。

\section{结论}
% 总结主要发现，说明局限与后续工作。

\bibliographystyle{gbt7714-numerical}
\bibliography{references}

\end{document}
